\clearpage\ArithmeticHead{425}{R}{KONVERGENZ DER ENTWICKLUNG DER MITTELPUNKTSGLEICHUNG.}
\begingroup\makeatletter\let\veqno\@@leqno\makeatother
\newcommand{\vsroot}[1]{\mathord{\vcenter{\hbox{\Large\(\surd\)}}}#1}
\begin{center}[4.]\end{center}
\noindent{\small [Dritter Zettel]}\par
\begin{center}\textbf{Nro. 2 als weiterer Zusatz.}\end{center}
\noindent\hspace*{\parindent}Indem man
\begin{equation}\begin{gathered}
\cos\varphi=p,\ \cos\theta=q,\\
{}[x=1-g\cos M,\ y=1-f\cos E]
\end{gathered}\tag*{[30]}\end{equation}
\noindent setzt, wird[\textsuperscript{*})]
\begin{equation}\left\{\begin{aligned}
\frac{px}{yy}={}&\frac12q\left(\begin{aligned}
 &1-\frac13\,{.}\,\frac1p\,{.}\,i\varepsilon-\frac1{12}\varepsilon\varepsilon+\frac1{60}\frac1p i\varepsilon^3\,{.}\,{.}\,{.}\\
 &\qquad+\frac14\,{.}\,\frac pq\varepsilon\varepsilon-\frac16\,{.}\,\frac1q i\varepsilon^3\,{.}\,{.}\,{.}
\end{aligned}\right)\\
 &\ \times\left(\begin{aligned}
 &1+\frac1p\,{.}\,i\varepsilon+\frac13\varepsilon\varepsilon+\frac5{12}\,{.}\,\frac1p i\varepsilon^3\,{.}\,{.}\,{.}\\
 &\qquad-\frac34\,{.}\,\frac1{pp}\varepsilon\varepsilon-\frac12\,{.}\,\frac1{p^3}i\varepsilon^3\,{.}\,{.}\,{.}
\end{aligned}\right)
\end{aligned}\right.\tag*{[31]}\end{equation}
\[\begin{aligned}
\left[\frac2q\,{.}\,\frac{px}{yy}=\right]1
 &+\frac23\,{.}\,\frac1p i\varepsilon+\frac14\varepsilon\varepsilon+\frac{13}{30}\,{.}\,\frac1p i\varepsilon^3\,{.}\,{.}\,{.}\\
 &+\frac14\,{.}\,\frac pq\varepsilon\varepsilon-\frac16\,{.}\,\frac1q i\varepsilon^3\,{.}\,{.}\,{.}\\
 &-\frac5{12}\,{.}\,\frac1{pp}\varepsilon\varepsilon-\frac12\,{.}\,\frac1{p^3}i\varepsilon^3\,{.}\,{.}\,{.}\\
 &\qquad-\frac7{36}\,{.}\,\frac1p i\varepsilon^3\,{.}\,{.}\,{.}\\
 &\qquad+\frac14\,{.}\,\frac1q i\varepsilon^3\,{.}\,{.}\,{.}\\
 &\qquad+\frac14\,{.}\,\frac1{p^3}i\varepsilon^3\,{.}\,{.}\,{.}
\end{aligned}\]
\par\vfill\noindent\rule{.25\linewidth}{.4pt}\par\smallskip
{\footnotesize\noindent\hspace*{\parindent}[\textsuperscript{*}) Die in der Handschrift an dieser Stelle folgende Gleichung ist unvollständig, die vollständige Gleichung findet sich auf der zweiten Seite des Blättchens; wir haben gleich diese hierher gesetzt und die unvollständige unterdrückt.]\par}
\noindent{\scriptsize X\,1.}\hfill{\small 54}\par
\clearpage\SourceLeaf{433}\ArithmeticHead{426}{L}{ANALYSIS. NACHLASS.}
\begin{equation}\begin{aligned}
\left[\frac2q\,{.}\,\frac{px}{yy}\right]={}&1+\frac23\frac1p i\varepsilon+\frac1{12ppq}(3ppq+3p^3-5q)\varepsilon\varepsilon\\
 &+\frac1{180p^3q}(43ppq+15p^3-45q)i\varepsilon^3\,{.}\,{.}\,{.}
\end{aligned}\tag*{[32]}\end{equation}
\begin{equation}\begin{aligned}
\varepsilon={}&-i\vsroot{\frac{2x}{pq}}-\frac13 i\frac{x}{ppq}\\
 &-\frac1{36}i\vsroot{\frac{2x^3}{pq}}\left\{\frac5{p^3q}-\frac3{pq}+\frac9{qq}\right\}\,{.}\,{.}\,{.}
\end{aligned}\tag*{[33]}\end{equation}
\[\begin{aligned}
\left[\frac2q\,{.}\,\frac{px}{yy}=\right]1
 &+\frac23\vsroot{\frac{2x}{p^3q}}+\frac29\,{.}\,\frac{x}{p^3q}\\
 &\qquad+\frac1{54}\vsroot{\frac{2x^3}{p^9q^5}}(5q-3ppq+9p^3)\\
 &-\frac16\,{.}\,\frac{x}{p^3qq}(3ppq+3p^3-5q)\\
 &\qquad-\frac1{18}\vsroot{\frac{2x^3}{p^9q^5}}(3ppq+3p^3-5q)\\
 &\qquad-\frac1{90}\vsroot{\frac{2x^3}{p^9q^5}}(43ppq+15p^3-45q)\\
 &\qquad{.}\quad{.}\quad{.}\quad{.}\quad{.}\quad{.}\quad{.}\quad{.}\quad{.}\quad{.}\quad{.}
\end{aligned}\]
\begin{equation}\begin{aligned}
={}&1+\frac23\vsroot{\frac{2x}{p^3q}}+\frac1{18}\,{.}\,\frac{x}{p^3qq}(19q-9ppq-9p^3)\\
 &+\frac1{270}\vsroot{\frac{2x^3}{p^9q^5}}(235q-189ppq-45p^3)\\
 &\quad{.}\quad{.}\quad{.}\quad{.}\quad{.}\quad{.}\quad{.}\quad{.}\quad{.}\quad{.}\quad{.}\quad{.}
\end{aligned}\tag*{[34]}\end{equation}
\noindent\hspace*{\parindent}Logarithm von
\begin{equation}\frac1{540}\vsroot{\frac2{p^9q^3}}(235q-189ppq-45p^3)\tag*{[35]}\end{equation}
\noindent ist \(9{,}9046115\). Die Multiplication mit \(\surd(1-g\cos M)\) auf der dritten Seite[\textsuperscript{*})].
\begin{center}\rule{.24\linewidth}{.4pt}\end{center}
\noindent\hspace*{\parindent}Näherungsweise wird \((n)\) dargestellt durch
\[\frac{G(\operatorname{tang}\frac12\theta)^n}{(n+2{,}6)^{\frac52}},\quad \log G=7{,}18800\]
\noindent oder auch
\[\frac{G(\operatorname{tang}\frac12\theta)^n}{(n+2{,}54)^{\frac52}},\quad [\log]G\ [=]\ 7{,}18234.\]
\begin{center}\rule{.24\linewidth}{.4pt}\end{center}
\par\vfill\noindent\rule{.25\linewidth}{.4pt}\par\smallskip
{\footnotesize\noindent\hspace*{\parindent}[\textsuperscript{*}) Die dritte und vierte Seite unseres Blättchens enthalten fast nur Zahlenreihen, die nicht abgedruckt sind.]\par}
\clearpage\endgroup
